\par\Needspace{6\baselineskip}\medskip
\begin{center}\hypertarget{native-236-article171}{}\pdfbookmark[2]{171.}{article-171}171.\end{center}
Exemplum ad illustrationem huius disquisitionis a \emph{Vesta} desumere conveniet, quae inter omnes planetas recentissime detectos inclinatione ad \SourceLeaf{237}eclipticam minima gaudet\footnote{Nihilominus haec inclinatio etiamnum satis considerabilis est, ut orbitae determinationem satis tuto atque exacte \emph{tribus} observationibus superstruere liceat: revera elementa prima, quae hoc modo ex observationibus 19 tantum diebus ab invicem distantibus deducta erant (vid. \textsc{von Zach} Monatl. Corresp. Vol. XV. p. 595), proxime iam accedunt ad ea, quae hic ex observationibus quatuor, 162 diebus ab invicem dissitis, derivabuntur.}. Eligimus observationes sequentes Bremae, Parisiis, Lilienthalii et Mediolani ab astronomis clarr. \textsc{Olbers}, \textsc{Bouvard}, \textsc{Bessel} et \textsc{Oriani} institutas:
\begin{center}
\begingroup\fontsize{9}{11}\selectfont\setlength{\tabcolsep}{3pt}
\begin{tabular}{|l|c|l|}
\hline
\shortstack{Tempus med. loci\\observationis} & Ascensio recta & Declinatio\\\hline
1807 Martii 30. $12^{\mathrm h}33^{\mathrm m}17^{\mathrm s}$ & $183^0\,52'\,40''{,}8$ & $11^0\,54'\,27''$ Bor.\\
\phantom{1807 }Maii 17. $8\quad16\quad5$ & $178\quad36\quad42{,}3$ & $11\quad39\quad46{,}8$ —\\
\phantom{1807 }Iulii 11. $10\quad30\quad19$ & $189\quad49\quad7{,}7$ & $3\quad9\quad10{,}1$ —\\
\phantom{1807 }Sept. 8. $7\quad22\quad16$ & $212\quad50\quad3{,}4$ & $8\quad38\quad17{,}0$ Austr.\\\hline
\end{tabular}\endgroup
\end{center}
Pro iisdem temporibus e tabulis motuum Solis invenimus
\begin{center}
\begingroup\fontsize{9}{11}\selectfont\setlength{\tabcolsep}{2pt}
\begin{tabular}{|l|c|c|c|c|c|}
\hline
& \shortstack{Longit. Solis ab\\aequin. app.} & Nutatio & \shortstack{Distantia a\\terra} & \shortstack{Latitudo\\Solis} & \shortstack{Obliquitas eclipt.\\apparens}\\\hline
Martii 30 & $9^0\,21'\,59''{,}5$ & $+16''{,}8$ & $0{,}999\,6448$ & $+0''{,}23$ & $23^0\,27'\,50''{,}82$\\
Maii 17 & $55\quad56\quad20{,}0$ & $+16{,}2$ & $1{,}0119789$ & $-0{,}63$ & $49{,}83$\\
Iulii 11 & $108\quad34\quad53{,}3$ & $+17{,}3$ & $1{,}0165795$ & $-0{,}46$ & $49{,}19$\\
Sept. 8 & $165\quad8\quad57{,}1$ & $+16{,}7$ & $1{,}0067421$ & $+0{,}29$ & $49{,}26$\\\hline
\end{tabular}\endgroup
\end{center}

Iam loca observata planetae, adhibita eclipticae obliquitate apparente, in longitudines et latitudines conversa, a nutatione et aberratione fixarum purgata, tandemque demta praecessione ad initium anni 1807 reducta sunt, dein e locis Solis ad normam praeceptorum art. 72 derivata sunt loca terrae ficta (ut parallaxis ratio habeatur), longitudinesque demta nutatione et praecessione ad eandem epocham translatae; tandem tempora ab initio anni numerata et ad meridianum Parisinum reducta. Hoc modo orti sunt numeri sequentes:
\begin{center}
\begingroup\fontsize{9}{11}\selectfont\setlength{\tabcolsep}{2.5pt}
\begin{tabular}{l|c|c|c|c}
$t,\ t',\ t'',\ t'''$ & $89{,}505\,162$ & $137{,}344502$ & $192{,}419502$ & $251{,}288102$\\
$a,\ a',\ a'',\ a'''$ & $178^0\,43'\,38''{,}87$ & $174^0\,1'\,30''{,}08$ & $187^0\,45'\,42''{,}23$ & $213^0\,34'\,15''{,}63$\\
$\beta,\ \beta',\ \beta'',\ \beta'''$ & $12\quad27\quad6{,}16$ & $10\quad8\quad7{,}80$ & $6\quad47\quad25{,}51$ & $4\quad20\quad21{,}63$\\
$l,\ l',\ l'',\ l'''$ & $189\quad21\quad33{,}71$ & $235\quad56\quad0{,}63$ & $288\quad35\quad20{,}32$ & $345\quad9\quad18{,}69$\\
$\mathrm{logg}\,R,R',R'',R'''$ & $9{,}9997990$ & $0{,}005\,1376$ & $0{,}0071739$ & $0{,}0030625$
\end{tabular}\endgroup
\end{center}

\SourceLeaf{238}Hinc deducimus
\[
\begin{array}{ll}
\gamma' = 168^{0}32'41''{,}34, & \delta' = 62^{0}23'4''{,}88 \\
\gamma'' = 173\ \ 5\ 15{,}68, & \delta'' = 100\ 45\ 1{,}40 \\
\log a' = 9{,}952\,6104, & b' = -11{,}009\,449 \\
\log a'' = 9{,}999\,4839, & b'' = -\ 2{,}082\,036
\end{array}
\]
\[
\begin{array}{ll}
\varkappa' = -1{,}083\,306, & \log \lambda = 0{,}072\,8800,\\
\log \mu' = 9{,}713\,9702_n & \\
\varkappa'' = +6{,}322\,006, & \log \lambda''' = 0{,}079\,8512_n,\\
\log \mu'' = 9{,}838\,7061 & 
\end{array}
\]
\[
\begin{array}{ll}
A'D = 37^{0}17'51''{,}50, & A''D = 89^{0}24'11''{,}84,\\
\epsilon = 9^{0}5'5''{,}48 & \\
B'D = -25\ \ 5\ 13{,}38, & B''D = -11\ 20\ 49{,}56.
\end{array}
\]
His calculis praeliminaribus absolutis, \emph{hypothesin primam} aggredimur. E
temporum intervallis elicimus
\[
\begin{array}{l}
\log k(t'-t) = 9{,}915\,3666 \\
\log k(t''-t') = 9{,}976\,5359 \\
\log k(t'''-t'') = 0{,}005\,4651
\end{array}
\]
atque hinc valores primos approximatos
\[
\begin{array}{ll}
\log P' = 0{,}06117, & \log(1+P') = 0{,}33269,\\
\log Q' = 9{,}59087 & \\
\log P'' = 9{,}97107, & \log(1+P'') = 0{,}28681,\\
\log Q'' = 9{,}68097, & 
\end{array}
\]
hinc porro
\[
\begin{array}{ll}
c' = -7{,}68361, & \log d' = 0{,}04666_n \\
c'' = +2{,}20771, & \log d'' = 0{,}12552.
\end{array}
\]
Hisce valoribus, paucis tentaminibus factis, solutio sequens aequationum I. II
elicitur:
\[
\begin{array}{ll}
x' = 2{,}04856, & z' = 23^{0}38'17'',\\
\log r' = 0{,}34951 & \\
x'' = 1{,}95745, & z'' = 27\ \ 2\ \ 0,\\
\log r'' = 0{,}34194. & 
\end{array}
\]
Ex $z'$, $z''$ atque $\epsilon$ eruimus $C'C'' = v''-v' = 17^{0}7'5''$: hinc
$v'-v$, $r$, $v'''-v''$, $r'''$ per aequationes sequentes determinandae erunt:
\[
\begin{array}{l}
\log r\ \sin(v'-v) = 9{,}74962, \\ \log r\ \sin(v'-v+17^{0}7'5'') = 0{,}07500 \\
\log r'''\sin(v'''-v'') = 9{,}84729, \\ \log r'''\sin(v'''-v'+17^{0}7'5'') = 0{,}10733,
\end{array}
\]
unde eruimus
\SourceLeaf{239}
\[
\begin{array}{ll}
v'-v = 14^{0}14'32'', & \log r = 0{,}35865 \\
v'''-v'' = 18\ 48\ 33, & \log r''' = 0{,}33887.
\end{array}
\]
Denique invenitur
\[
\begin{gathered}
\log(\eta 01) = 0{,}00426, \\ \log(\eta 12) = 0{,}00599, \\ \log(\eta 23) = 0{,}00711,
\end{gathered}
\]
atque hinc valores correcti ipsarum $P'$, $P''$, $Q'$, $Q''$:
\[
\begin{array}{ll}
\log P' = 0{,}05944, & \log Q' = 9{,}60374 \\
\log P'' = 9{,}97219, & \log Q'' = 9{,}69581,
\end{array}
\]
quibus \emph{hypothesis secunda} superstruenda erit. Huius praecipua momenta
ita se habent:
\[
\begin{array}{ll}
c' = -7{,}67820, & \log d' = 0{,}045\,736_n\mathord{\rule[.1ex]{.35em}{.35pt}} \\
c'' = +2{,}21061, & \log d'' = 0{,}126\,054
\end{array}
\]
\[
\begin{array}{ll}
x' = 2{,}03308, & z' = 23^{0}47'54'',\\
\log r' = 0{,}346\,747 & \\
x'' = 1{,}94290, & z'' = 27\ 12\ 25,\\
\log r'' = 0{,}339\,373 & 
\end{array}
\]
\[
C'C'' = v''-v' = 17^{0}8'0''
\]
\[
\begin{array}{ll}
v'-v = 14^{0}21'36'', & \log r = 0{,}354\,687 \\
v'''-v'' = 18\ 50\ 43, & \log r''' = 0{,}334\,564
\end{array}
\]
\[
\begin{gathered}
\log(\eta 01) = 0{,}004\,359, \\ \log(\eta 12) = 0{,}006\,102, \\ \log(\eta 23) = 0{,}007\,280.
\end{gathered}
\]
Hinc prodeunt valores denuo correcti ipsarum $P'$, $P''$, $Q'$, $Q''$:
\[
\begin{array}{ll}
\log P' = 0{,}059\,426, & \log Q' = 9{,}604\,749 \\
\log P'' = 9{,}972\,249, & \log Q'' = 9{,}697\,564,
\end{array}
\]
quibus si ad \emph{tertiam hypothesin} progredimur, numeri sequentes
resultant:
\[
\begin{array}{ll}
c' = -7{,}67815, & \log d' = 0{,}045\,729_n \\
c'' = +2{,}21076, & \log d'' = 0{,}126\,082
\end{array}
\]
\[
\begin{array}{ll}
x' = 2{,}03255, & z' = 23^{0}48'14'',\\
\log r' = 0{,}346\,653 & \\
x'' = 1{,}94235, & z'' = 27\ 12\ 49,\\
\log r'' = 0{,}339\,276 & 
\end{array}
\]
\[
C'C'' = v''-v' = 17^{0}8'4''
\]
\[
\begin{array}{ll}
v'-v = 14^{0}21'49'', & \log r = 0{,}354\,522 \\
v'''-v'' = 18\ 51\ \ 7, & \log r''' = 0{,}334\,290
\end{array}
\]
\[
\begin{gathered}
\log(\eta 01) = 0{,}004\,363, \\ \log(\eta 12) = 0{,}006\,106, \\ \log(\eta 23) = 0{,}007\,290.
\end{gathered}
\]
\SourceLeaf{240}
Quodsi iam ad normam praeceptorum art.\ praec.\ distantiae a terra
supputantur, prodit:
\[
\begin{array}{ll}
\rho' = 1{,}5635, & \rho'' = 2{,}1319 \\
\log \rho \cos\beta = 0{,}09876, & \log \rho'''\cos\beta''' = 0{,}42842 \\
\log \rho \sin\beta = 9{,}44252, & \log \rho'''\sin\beta''' = 9{,}30905 \\
\beta = 12^{0}26'40'', & \beta''' = 4^{0}20'39'' \\
\log \rho = 0{,}10909, & \log \rho''' = 0{,}42967.
\end{array}
\]
Hinc inveniuntur
\[
\begin{array}{c|c|c}
 & \text{Correctiones temporum} & \text{Tempora correcta} \\
\hline
\text{I} & 0{,}007\,335 & 89{,}497\,827 \\
\text{II} & 0{,}008\,921 & 137{,}335\,581 \\
\text{III} & 0{,}012\,165 & 192{,}407\,337 \\
\text{IV} & 0{,}015\,346 & 251{,}272\,756
\end{array}
\]
unde prodeunt valores quantitatum $P'$, $P''$, $Q'$, $Q''$ denuo correcti
\[
\begin{array}{ll}
\log P' = 0{,}059\,415, & \log Q' = 9{,}604\,782 \\
\log P'' = 9{,}972\,253, & \log Q'' = 9{,}697\,687.
\end{array}
\]
Tandem si hisce valoribus novis \emph{hypothesis quarta} formatur, numeri
sequentes prodeunt:
\[
\begin{array}{ll}
c' = -7{,}678\,116, & \log d' = 0{,}045\,723_n \\
c'' = +2{,}210\,773, & \log d'' = 0{,}126\,084
\end{array}
\]
\[
\begin{array}{ll}
x' = 2{,}032\,473, & z' = 23^{0}48'16''{,}7,\\
\log r' = 0{,}346\,638 & \\
x'' = 1{,}942\,281, & z'' = 27\ 12\ 51{,}7,\\
\log r'' = 0{,}339\,263 & 
\end{array}
\]
\[
\begin{gathered}
v''-v' = 17^{0}8'5''{,}1, \\ \tfrac{1}{2}(u''+u') = 176^{0}7'50''{,}5, \\
\tfrac{1}{2}(u''-u') = 4^{0}33'23''{,}6
\end{gathered}
\]
\[
\begin{array}{ll}
v'-v = 14^{0}21'51''{,}9, & \log r = 0{,}354\,503 \\
v'''-v'' = 18\ 51\ \ 9{,}5, & \log r''' = 0{,}334\,263.
\end{array}
\]
Hi numeri ab iis, quos hypothesis tertia suppeditaverat, tam parum
differunt, ut iam tuto ad ipsorum elementorum determinationem progredi
liceat. Primo situm plani orbitae eruimus. Per praecepta art.\ 149
invenitur, ex $\gamma'$, $u'$ atque $A'C' = \delta'-z'$, inclinatio orbitae
$= 7^{0}8'14''{,}8$, longitudo nodi ascendentis $103^{0}16'37''{,}2$, argumentum
latitudinis in observatione secunda $94^{0}36'4''{,}9$, \SourceLeaf{241}adeoque longitudo in orbita $197^{0}52'42''{,}1$; perinde ex $\gamma''$, $u''$
atque $A''C'' = \delta''-z''$ elicitur inclinatio orbitae $= 7^{0}8'14''{,}8$,
longitudo nodi ascendentis $103^{0}16'37''{,}5$, argumentum latitudinis in
observatione tertia $111^{0}44'9''{,}7$, adeoque longitudo in orbita
$215^{0}0'47''{,}2$. Hinc erit longitudo in orbita pro observatione prima
$183^{0}30'50''{,}2$, pro quarta $233^{0}51'56''{,}7$. Quodsi iam ex $t'''-t$,
$r$, $r'''$ atque $v'''-v = 50^{0}21'6''{,}5$ orbitae dimensiones
determinantur, prodit:
\begin{center}
\begingroup\fontsize{9}{11}\selectfont\setlength{\tabcolsep}{3pt}
\begin{tabular}{@{}lr@{}}
Anomalia vera pro loco primo . . . . . . . . . . . . . . . . . & $293^{0}33'43''{,}7$ \\
Anomalia vera pro loco quarto . . . . . . . . . . . . . . . . . & $343\ \ 54\ \ 50{,}2$ \\
Hinc longitudo perihelii . . . . . . . . . . . . . . . . . . . . . . . & $249\ \ 57\ \ \ 6{,}5$ \\
Anomalia media pro loco primo . . . . . . . . . . . . . . . . & $302\ \ 33\ \ 32{,}6$ \\
Anomalia media pro loco quarto . . . . . . . . . . . . . . . & $346\ \ 32\ \ 25{,}2$ \\
Motus medius diurnus sidereus . . . . . . . . . . . . . . . . . & $978''{,}7216$ \\
Anomalia media pro initio anni 1807 . . . . . . . . . . . & $278^{0}13'39''{,}1$ \\
Longitudo media pro eadem epocha . . . . . . . . . . . . & $168\ \ 10\ \ 45{,}6$ \\
Angulus $\varphi$ . . . . . . . . . . . . . . . . . . . . . . . . . . . . . . . . . . & $5\ \ \ 2\ \ 58{,}1$ \\
Logarithmus semiaxis maioris . . . . . . . . . . . . . . . . . . & $0{,}372\,898.$
\end{tabular}\endgroup
\end{center}

Si secundum haecce elementa pro temporibus $t$, $t'$, $t''$, $t'''$ correctis
loca planetae geocentrica computantur, quatuor longitudines cum $a$, $a'$,
$a''$, $a'''$, duaeque latitudines intermediae cum $\beta'$, $\beta''$ ad
unam minuti secundi partem decimam conspirant; latitudines extremae vero
prodeunt $12^{0}26'43''{,}7$ atque $4^{0}20'40''{,}1$, illa $22''{,}4$ errans
defectu, haec $18''{,}5$ excessu. Attamen si manentibus elementis reliquis
tantummodo inclinatio orbitae $6''$ augeatur, longitudoque nodi $4' 40''$
diminuatur, errores inter omnes latitudines distributi ad pauca minuta
secunda deprimentur, longitudinesque levissimis tantum erroribus
afficientur, qui et ipsi propemodum ad nihilum reducentur, si insuper
epocha longitudinis $2''$ diminuatur.
\par\medskip\begin{center}\rule{.24\linewidth}{.35pt}\end{center}
